\documentclass{article}
\usepackage{amsmath,amssymb,amsthm}
\title{7.4: Integration of Rational Functions by Partial Fractions}
\author{Amarnath Patel}
\date{\today}

\begin{document}

\maketitle

$f(x) = \frac{P(x)}{Q(x)}$ 
\newline
Where P and Q are polynomials and the degree of P is less than the degree of Q.\\
Its possible to express f as a sum of simpler fractions provided that the degree of P is less than teh degree of Q. These are called proper.\\
If f is improper that is, the degree of P is greater than Q, we must take the preliminary step of dividing Q into P via long division, until a remainder R(x) is resolved, such that the degree of R is less than Q.
This results in the expression:
\newline
\centering
$f(x) = \frac{P(x)}{Q(x)} = S(x) + \frac{R(x)}{Q(x)}$
\newline
\flushleft
Where S(x) is a polynomial and the degree of R is less than the degree of Q.\\

If the denominator Q(x) is factorable the next step is to factor out Q(x) as far as possible. 
The third step is to express proper rational function $R(x)/Q(x)$ as a sum of \textbf{partial fractions} of the form:
$\frac{A}{ax+b}^i$ or $\frac{Ax+B}{(ax^2+bx+c)^j}$

The first case is where Q(x) is a product of distinct linear factors.
Q(x) = $(ax+b)^i_1 (ax+b)^i_2 ... (ax+b)^i_n$ where $i_1, i_2, ... i_n$ are positive integers.\\

$\frac{R(x)}{Q(x)} = \frac{A_1}{ax+b} + \frac{A_2}{(ax+b)^2} + ... + \frac{A_n}{(ax+b)^n}$

$\int \frac{x^2 +2x -1}{2x^3 +3x^2 -2x}dx$

$2x^3 +3x^2 -2x = x(2x^2 +3x -2) = x(2x-1)(x+2)$

$\frac{x^2 +2x-1 }{x(2x-1)(x+2)} = \frac{A}{x} + \frac{B}{2x-1} + \frac{C}{x+2}$

\section{Example 2}

$\int \frac{dx}{x^2-a^2}$ where $a \neq 0$

$\frac{1}{x^2-a^2} = \frac{1}{(x-a)(x+a)} = \frac{A}{x-a} + \frac{B}{x+a}$

Therefore:

$A(x+a) + B(x-a) = 1$

$\frac{1}{2a} \int (\frac{1}{x-a} - \frac{1}{x+a})dx = \frac{1}{2a}(\ln|x-a| - \ln|x+a|) + C$

\section{Example 3 and case 2}

Q(x) is a product of linear factors, some of which are repeated.

$\int \frac{x^4 -2x^2 +4x+ 1}{x^3 +x^2 -x +1}dx$
$x+1 + \frac{4x}{x^3 +x^2 -x +1}$
Next we need to factor out the denominator Q(x) as far as possible. Since Q(1) = 0 we know that x-1 is a factor.
$(x-1)^2(x+1)$
Since the linear factor x-1 occurs twice, the partial fraction decomposition is 
$\frac{A}{x-1} + \frac{B}{(x-1)^2} + \frac{C}{x+1}$
$4x = A(x-1)(x+1)+B(x+1)+C(x-1)^2$
$ = (A+C)x^2 + (B-2C)x + (-A +B + C)$
$A + C = 0$
$B-2C = 4$
$-A + B + C = 0$
A = 1, B = 2, C = -1, so
\newline
$\frac{x^2}{2} + x - \frac{2}{x-1} + \ln|\frac{x-1}{x+1}| + K$

\section{Case 3}
Q(x)contains irreducible quadratic factors, none of which is repeated.

$\int \frac{2x^2-x+4}{x^3+4x}dx$
\newline
$\frac{2x^2-x+4}{x^3+4x} = \frac{A}{x} + \frac{Bx+C}{x^2+4}$
\newline
$2x^2-x+4 = A(x^2+4) + (Bx+C)x$
\newline
$= (A+B)x^2 + Cx + 4A$
\newline
$A+B = 2$
$C = -1$
$4A = 4$
\newline
Therefore A and B = 1 and C = -1.
$u = x^2 + 4$
\newline
$du = 2xdx$
\newline
$\int \frac{1}{x}dx + \int \frac{x}{x^2+4}dx - \int \frac{1}{x^2+4}dx$
\newline
$ = \ln|x| + \frac{1}{2} \ln|x^2+4| - \frac{1}{2} \tan^{-1}(\frac{x}{2}) + K$
\newline

\section{Case 4}

Q(x) contains a repeated irreducible quadratic factor.

$\int \frac{1-x+2x^2-x^3}{x(x^2 +1)^2}dx$
\newline
$ = \frac{A}{x} + \frac{Bx+C}{x^2+1} + \frac{Dx+E}{(x^2+1)^2}$
\newline
$-x^3 + 2x^2 -x + 1 = A(x^2+1)^2 + (Bx+C)x(x^2+1) + (Dx+E)x$
\newline
$ = A(x^4 +2x^2+1) + B(x^4 +x^2)+ C(x^3 +x) + Dx^2 + Ex$
\newline
$ = (A+B)x^4 + Cx^3 + (2A+B+D)x^2 + (C+E)x + A$

$A+B = 0$
$C = -1$
$2A+B+D = 2$
$C+E = -1$
$A = 1$
Each therefore:
$B = -1$
$D = 1$
$E = 0$
$A = 1$
$C = -1$

$\int (\frac{1}{x} - \frac{x+1}{x^2+1} + \frac{x}{(x^2+1)^2})dx$
$ln|x|-\frac{1}{2}ln(x^2+1)-\tanh(x)-\frac{1}{2(x^2+1)} + K$


\section{Rationalizing Substitutions}

$\int \frac{\sqrt{x+4}}{x}dx$
\newline
Let $u = \sqrt{x+4}$. Then $x = u^2 -4$ and $dx = 2udu$.
\newline
$\int \frac{u}{u^2-4}2udu$
\newline
$\int (1+\frac{4}{u^2-4})du$
\newline
$u^2 -4 = (u-2)(u+2)$
\newline
$2\sqrt{x+4}+2\ln|\frac{\sqrt{x+4}-2}{\sqrt{x+4}+2}| + C$

\end{document}
